\section{Decoupling State from Execution}
\label{sec:approach}

\subsection{Storage Scalability and Security Trade-off}

\lijl{Outline:
\begin{itemize}
    \item Tension between storage scalability and strong security
    \item Full state replication guarantees security, but no storage scalability
    \item Sharding scales blockchain storage, but compromises security
    \item This trade-off seems to be fundamental:
    \item To provide strong security, \ie, the system can tolerate any $f$ adversaries, the number of nodes executing each transaction needs to be at least $f+1$.
    \item When adding more nodes to the system, $f$ also increases proportionally
    \item That means, proportionally more nodes need to maintain the state to execute the transactions
    \item The total state capacity of the system remains unchanged
    \item We can similarly show that scaling storage will result in weaker security
\end{itemize}
}

\subsection{The Case for Decoupling State from Execution}

\lijl{Outline:
\begin{itemize}
    \item This trade-off stems from a deep-rooted architectural decision
    \item In existing blockchain architectures, state and transaction execution are tightly coupled
    \item A node only executes a transaction (or part of a transaction) if it has all the state accessed in the transaction
    \item In traditional blockchain, each node maintains a full copy of the state, and executes all transactions
    \item In sharded blockchains, a node only executes transactions that access its responsible shards
    \item In stateless blockchains, only external nodes execute transactions
    \item Given this tight coupling, it is fundamentally hard to overcome the security and storage scalability trade-off
    \item In this work, we argue for a new blockchain architecture, where transaction execution and state storage are logically decoupled
    \item Clear separation of responsibility: blockchain execution performs transaction statements but is stateless; state storage only maintains blockchain world state, and is agnostic to transaction logic.
    \item This decoupling enables simultaneous storage scalability and the strongest security guarantees.
\end{itemize}
}

From prior research on fully replicated blockchains, sharded blockchains, and stateless blockchains, we observe a common principle: nodes store the state necessary to execute transactions independently.
In other words, the set of transactions a node can execute is determined by the state it maintains locally.
In fully replicated blockchains, every node maintains the complete state and executes all transactions.
In sharded blockchains, nodes only maintain a shard of the state, as they execute transactions that access that shard.
In stateless blockchains, internal nodes do not maintain state because they neither execute nor are able to execute transactions.
Instead, external storage nodes maintain all or part of the state and execute those transactions that can be supported by the state they store.

This ``store what you need for execution'' common theme leads to the dilemma between storage scalability and security guarantees.
On one hand, storage scalability cannot be achieved if the state is replicated to all nodes or a set of nodes whose size is proportional to the total number of nodes.
That means that for each transaction, the state it accesses can only be stored on a limited set of nodes whose number is independent of the total number of nodes.
If only the nodes storing the accessed state can execute the transaction, the number of executing nodes is also limited.
As the number of $f$ faulty nodes increases with the total number of nodes, the faulty nodes can eventually outnumber the executing nodes, making it impossible to achieve security guarantees (without seeking unconventional extensions to the security model).
This fundamental theoretical limitation forces sharded and stateless blockchains to trade off security guarantees in various ways (\autoref{sec:bg}) in exchange for storage scalability.
\lijl{This is an important insight. We can make this part even stronger.}
%\lijl{We can probably frame this into an informal impossibility result. That is, a system can choose at most one of the two properties. Then informally ``prove'' that if security is chosen, then storage scalability can't be achieved, and vice versa.}

\paragraph{State-execution decoupling.}
On the other hand, we observe that serving consistent state is a fundamentally easier problem than ensuring consistent execution.
The state can be self-verifiable, and receivers can be sure they obtain the correct state even if it is served by single data source.
This is unlike execution which is generally not verifiable without re-execution, so $f+1$ nodes need to execute the same transaction.
Based on this observation, we propose to \emph{decouple state from execution}.

With the proposed approach, each transaction is executed by \emph{all} nodes, regardless of the state they maintain.
The nodes are allowed to maintain only a portion of the state or even no state at all.
We exploit the fact that while the entire state is necessary to execute all possible transactions, each \emph{individual} transaction may not access the entire state, and most transactions in practice access only a small portion of the state.
When executing a transaction, if a node needs to access state that it does not maintain locally, it retrieves the required state from other nodes to complete the execution.
With this state-execution decoupling, we can independently design state management to achieve storage scalability, while all nodes execute all transactions to retain security guarantees.
\lijl{Can relate to other systems designs, e.g., memory is separated from CPU execution logic (CPU only maintains registers and caches), prior disaggregated memory.}
If a transaction accesses a state shard not stored locally, the node retrieves it from other nodes to complete the execution.
By carefully designing the placement of state shards, we ensure that each shard is stored on a limited subset of nodes, thereby achieving storage scalability.
\lijl{A key principle we are arguing here: strong security requires nodes proportional to $f$ to execute the transactions, but availability of data can be done in a scalable fashion.}

\lijl{Another point to make: why do nodes need to maintain some state locally, and how much to store? There are two reasons: state availability and performance. State availability determines the storage lower bound on each node. Storing more state than the lower bound could improve performance.}

\paragraph{The size of the retrieved temporary state.}
A natural concern with this design is that, in the worst case, the size of the retrieved temporary state could be as large as the entire state, for instance, if the transactions scan the entire state.
We emphasize that this is a theoretical worst case and is unlikely to occur in practice.
Firstly, as our design will demonstrate into \autoref{sec:overview}, nodes retrieve individual key-value pairs from each other instead of transferring the whole state portion the remote nodes maintain.
Thus, the peak size of the retrieved temporary state equals to the \emph{working set}, i.e., the total size of the accessed keys and values, of the executing transactions.
Empirically, most transactions access only an extremely small fraction of the state.
According to the data over 100 days in 2024 from Ethereum, on average each block reads ~1600 slots, about 0.0001\% of the total 151 million slots in the state~\cite{ProperDiskIO_GasPricing_LRUcache_2024}.
In practice, blockchain nodes preserve the working set of one or more recent blocks entirely in memory for performance reasons~\cite{nodereal2022_bsc_cache_hit_rate}.
We follow this practice in our design, assuming trivial memory footprint for the retrieved temporary state without loss of generality.
\lijl{Yes, should provide more concrete evidence and workload citations. Even better if we can do some empirical analysis of available traces.} \gd{Kind of elaborated.}

% However, diverging from the common practice also sacrifices its advantages: executing with local state is efficient and does not require additional security analysis.
% When the state is maintained locally, it is always correct and available to the node, and executing the same sequence of transactions from the same initial state will deterministically yield the same final state and output.
% Consequently, prior sharded and stateless blockchains can primarily focus on establishing consistent transaction ordering to guarantee their security properties.
% In contrast, our approach of decoupling state from execution introduces new challenges in ensuring the correctness and availability of state shards retrieved from other nodes, which are essential for security guarantees.
% If a node retrieves incorrect values, execution may diverge.
% If retrieval fails to complete, execution is stalled.
% Even if retrieval completes but is inefficient, execution will be delayed and throughput will decrease.
% We will elaborate on these challenges and outline our solutions in \autoref{sec:overview}; before that, we first formalize the interfaces and security properties of state management.
% \lijl{I think we can move these challenge discussions later.}
